% Part:sets-functions-relations
% Chapter: size-of-sets
% Section: non-enumerability-alt

\documentclass[../../../include/open-logic-section]{subfiles}

\begin{document}

\olfileid{sfr}{siz}{nen-alt}

\olsection{\printtoken{S}{nonenumerable} সেট}

\begin{editorial}
  এই অংশে \olref[enm-alt]{sec}-এর সংজ্ঞাগুলি ব্যবহার করে $\Bin^\omega$ ও $\Pow{\Nat}$ যে অতালিকায়নযোগ্য তা প্রমাণ করা হয়েছে; অর্থাৎ, $\PosInt$ থেকে সমাপতিত অপেক্ষকের বদলে~$\Nat$-এর সঙ্গে একৈক সমাপতন চাওয়া হয়েছে।
\end{editorial}

\begin{explain}
স্বাভাবিক সংখ্যার সেট $\Nat$ অসীম। আবার সেটি স্পষ্টতই !!{enumerable}। কিন্তু বিস্ময়কর সত্যটি হলো,
\emph{!!{nonenumerable}} সেটও আছে, অর্থাৎ এমন সেট যা !!{enumerable} নয়
(\olref[sfr][siz][enm-alt]{defn:enumerable} দেখো)।

এ কথা বিস্ময়কর মনে হতে পারে। কারণ $A$ !!{nonenumerable} বলার অর্থ হলো,
!!{bijection} $f \colon \Nat \to A$ একেবারেই \emph{নেই}; অর্থাৎ,~$\Nat$-এর অসীমসংখ্যক !!{element}s-কে~$A$-তে পাঠানো কোনো অপেক্ষকই~$A$-এর সমস্ত উপাদান অন্তর্ভুক্ত করে না। তাই $A$ !!{nonenumerable} হলে, স্বাভাবিক সংখ্যা যতগুলি আছে তার চেয়ে~$A$-এর !!{element}s ``বেশি''।

কোনো সেট !!{nonenumerable} প্রমাণ করতে হলে দেখাতে হবে যে উপযুক্ত কোনো !!{bijection} থাকতে পারে না। এর সবচেয়ে ভালো উপায় হলো দেখানো যে~$A$-এর !!{element}s তালিকায়িত করার প্রতিটি চেষ্টায় অন্তত একটি !!{element} বাদ পড়বেই; এতে প্রমাণ হয় যে কোনো অপেক্ষক $f\colon \Nat \to A$ !!{surjective} নয়। এটি প্রতিষ্ঠা করার একটি সাধারণ কৌশল হলো কান্টরের \emph{কর্ণ পদ্ধতি} ব্যবহার করা। $A$-এর !!{element}s-এর একটি তালিকা, ধরা যাক $x_1$, $x_2$, \dots, দেওয়া থাকলে আমরা~$A$-এর এমন আরেকটি !!{element} তৈরি করি, যা তার নির্মাণের কারণেই ওই তালিকায় থাকা অসম্ভব।

তবে একটি উদাহরণেই সব কথা সবচেয়ে ভালো বোঝা যায়। তাই আমাদের প্রথম উদাহরণ হলো $0$ ও $1$ দিয়ে গঠিত সমস্ত অসীম প্রতীকক্রমের সেট~$\Bin^\omega$। (`$\Bin$' দিয়ে বাইনারি বোঝানো হয়, এবং আমরা একে স্রেফ দুই-উপাদানের সেট
$\{0,1\}$ হিসেবে ভাবতে পারি।)\oliflabeldef{sfr:card-arithmetic:card-opps:sec}{\footnote{আরও নিখুঁতভাবে, আমাদের বলা উচিত যে $\Bin^\omega$ হলো $0$ ও $1$-এর সমস্ত $\omega$-অনুক্রমের সেট, অর্থাৎ সেটটি
$\funfromto{\omega}{\{0,1\}}$। কিন্তু \olref[sfr][card-arithmetic][card-opps]{sec}-এ গেলেই এর অর্থ স্পষ্ট হবে।}}{} আপাতত অবশ্য এই সামান্য শিথিল ভাষা কোনো বিভ্রান্তি সৃষ্টি করার কথা নয়।
\end{explain}

\begin{thm}
\ollabel{thm:nonenum-bin-omega}
$\Bin^\omega$~!!{nonenumerable}।
\end{thm}

\begin{proof}
$\Bin^\omega$-এর কোনো উপসেটের যেকোনো তালিকায়ন বিবেচনা করো। অর্থাৎ, আমাদের কাছে $s_{0}$, $s_{1}$, $s_{2}$, \dots{} এমন একটি তালিকা আছে, যেখানে প্রত্যেক $s_n$ হলো $0$ ও~$1$ দিয়ে গঠিত একটি অসীম প্রতীকক্রম। $s_n(m)$ দ্বারা এই তালিকার $n$-তম প্রতীকক্রমের $m$-তম অঙ্কটি বোঝানো হোক। % BN-SRC-003 TARGET-CORRECTION-BEGIN
\par\smallskip\noindent\textnormal{\small\textbf{উৎস-সংশোধন (BN-SRC-003):} হিমায়িত ইংরেজি উৎসে সূচকদুটির বর্ণনা উল্টে $s_n(m)$-কে ``$m$-তম প্রতীকক্রমের $n$-তম অঙ্ক'' বলা হয়েছে। পরের বাক্য, সারণি এবং প্রচলিত সংকেত অনুসারে এটি $n$-তম প্রতীকক্রমের $m$-তম অঙ্ক; বাংলা পাঠে সেই অর্থ দেওয়া হয়েছে।} % BN-SRC-003 TARGET-CORRECTION-END
তাহলে তালিকাটিকে এখন একটি সারণি হিসেবে ভাবতে পারি, যেখানে $s_n(m)$-কে $n$-তম সারি এবং $m$-তম স্তম্ভে বসানো হয়েছে:
\[
\begin{array}{c|c|c|c|c|c}
& 0 & 1 & 2 & 3 & \dots \\\hline
0 & \mathbf{s_{0}(0)} & s_{0}(1) & s_{0}(2) & s_0(3) & \dots \\\hline
1 & s_{1}(0)& \mathbf{s_{1}(1)} & s_1(2) & s_1(3) & \dots \\\hline
2 & s_{2}(0)& s_{2}(1) & \mathbf{s_2(2)} & s_2(3) & \dots \\\hline
3 & s_{3}(0)& s_{3}(1) & s_3(2) & \mathbf{s_3(3)} & \dots \\\hline
\vdots & \vdots & \vdots & \vdots & \vdots & \mathbf{\ddots}
\end{array}
\]
এবার আমরা $0$ ও $1$ দিয়ে গঠিত এমন একটি অসীম প্রতীকক্রম $d$ তৈরি করব, যা এই তালিকায় নেই। তার প্রত্যেকটি পদ নির্দিষ্ট করে, অর্থাৎ সকল $n \in \Nat$-এর জন্য $d(n)$ নির্দিষ্ট করে, আমরা এটি করব। স্বজ্ঞাতভাবে, উপরের সারণির কর্ণ বরাবর নেমে (তাই এর নাম ``কর্ণ পদ্ধতি'') প্রত্যেক $0$-কে $1$ এবং প্রত্যেক $1$-কে~$0$ করে দিই। % BN-SRC-004 TARGET-CORRECTION-BEGIN
\par\smallskip\noindent\textnormal{\small\textbf{উৎস-সংশোধন (BN-SRC-004):} হিমায়িত ইংরেজি উৎসের এই বাক্যে ``প্রত্যেক $1$-কে $0$ এবং প্রত্যেক $1$-কে $0$'' লেখা হয়েছে। সঙ্গে দেওয়া ক্ষেত্রভিত্তিক সংজ্ঞা অনুযায়ী প্রথম পরিবর্তনটি $0$ থেকে $1$; বাংলা পাঠে সেটি ঠিক করা হয়েছে।} % BN-SRC-004 TARGET-CORRECTION-END
আরও বিমূর্তভাবে, কর্ণের $n$-তম !!{element} $s_n(n)$ যথাক্রমে $1$ বা $0$ হওয়া অনুসারে $d(n)$-কে $0$ বা $1$ হিসেবে সংজ্ঞায়িত করি, অর্থাৎ:
\[
d(n) =
\begin{cases}
1 & \text{যদি $s_{n}(n) = 0$ হয়}\\
0 & \text{যদি $s_{n}(n) = 1$ হয়}
\end{cases}
\]
স্পষ্টতই $d \in \Bin^\omega$, কারণ এটি $0$ ও $1$ দিয়ে গঠিত একটি অসীম প্রতীকক্রম। কিন্তু আমরা $d$-কে এমনভাবে তৈরি করেছি যে যেকোনো $n \in \Nat$-এর জন্য $d(n) \neq s_n(n)$। অর্থাৎ, $d$ ও $s_n$ তাদের $n$-তম পদে আলাদা। তাই যেকোনো $n\in \Nat$-এর জন্য $d \neq s_n$। অতএব $d$ তালিকা $s_0$, $s_1$, $s_2$,
\dots{}-তে থাকতে পারে না।

$\Bin^\omega$-এর কোনো উপসেটের একটি নির্বিচার তালিকায়ন দেওয়া হলে সেটি যে~$\Bin^\omega$-এর কোনো !!{element} বাদ দেবেই, তা আমরা দেখিয়েছি। অতএব $\Bin^\omega$ সেটের কোনো তালিকায়ন নেই, অর্থাৎ $\Bin^\omega$ !!{nonenumerable}।
\end{proof}

\begin{explain}
এই প্রমাণপদ্ধতিকে ``কর্ণীকরণ'' বলা হয়, কারণ এখানে $d$ সংজ্ঞায়িত করতে সারণির কর্ণ ব্যবহার করা হয়েছে। তবে কর্ণীকরণে কোনো সারণি থাকা আবশ্যক নয়। বাস্তবে, কোনো সারণি ও প্রকৃত কর্ণ না থাকলেও একই ধরনের ধারণা ব্যবহার করে কোনো সেট যে
!!{nonenumerable} তা দেখানো যায়। পরের ফলটি এর একটি উদাহরণ।
\end{explain}

\begin{thm}
\ollabel{thm:nonenum-pownat}
$\Pow{\Nat}$ !!{enumerable} নয়।
\end{thm}

\begin{proof}
একই পথে এগিয়ে আমরা দেখাব যে~$\Nat$-এর উপসেটগুলির প্রত্যেকটি তালিকা $\Nat$-এর কোনো-না-কোনো উপসেট বাদ দেয়। তাই ধরি, $\Nat$-এর উপসেটগুলির কোনো তালিকা $N_0, N_1, N_2, \ldots$ আমাদের কাছে আছে। নিচের মতো করে একটি সেট $D$ সংজ্ঞায়িত করি: $n \in D$ যদি এবং কেবল যদি $n \notin N_{n}$:
\[
D = \Setabs{n \in \Nat}{n \notin N_n}
\]
স্পষ্টতই $D\subseteq \Nat$। কিন্তু $D$ তালিকায় থাকতে পারে না। কারণ নির্মাণ অনুসারে $n \in D$ যদি এবং কেবল যদি $n\notin N_n$, তাই যেকোনো $n \in \Nat$-এর জন্য $D \neq N_n$।
\end{proof}

\begin{explain}
আগের প্রমাণে কোনো কর্ণের উল্লেখ ছিল না। তবুও এইভাবে ছবি আঁকলে সেটির মধ্যে একটি কর্ণ আছে বলে ভাবতে পারো: কল্পনা করো, সেটগুলি $N_0$, $N_1$, \dots{} একটি সারণিতে লেখা হয়েছে; সেখানে $N_n$-কে $n$-তম সারিতে লিখতে $m \in N_n$ হলে এবং কেবল তখনই $m$-তম স্তম্ভে $m$ লিখি। যেমন, ধরা যাক তালিকার প্রথম চারটি সেট
$\{0,1,2,\dots\}$, $\{1, 3, 5, \dots\}$, $\{0,1,4\}$, এবং
$\{2,3,4,\dots\}$; তাহলে আমাদের সারণির শুরু হবে
\[
\begin{array}{r@{}rrrrrrr}
  N_0 = \{ & \mathbf{0}, & 1, & 2, & & & & \dots\}\\
  N_1 = \{ &  & \mathbf{1}, &  & 3, &  & 5, & \dots\}\\
  N_2 = \{ & 0, & 1, &  &  & 4\phantom{,} &  & \}\\
  N_3 = \{ &  &  & 2, & \mathbf{3}, & 4, & & \dots\}\\
  &\vdots & & & & & & \ddots\phantom{\}}
  \end{array}
\]
তারপর $D$ হলো কর্ণ বরাবর নেমে পাওয়া সেই সেট, যেখানে $n$ কর্ণে \emph{না} থাকলে এবং কেবল তখনই $n
\in D$ রাখি। তাই উপরের ক্ষেত্রে আমরা $0$ ও $1$ বাদ দেব,~$2$ অন্তর্ভুক্ত করব,~$3$ বাদ দেব, ইত্যাদি।
\end{explain}

\begin{prob}
একটি স্পষ্ট কর্ণ-যুক্তির সাহায্যে দেখাও যে সব অপেক্ষক $f \colon \Nat \to \Nat$-এর সেটটি
!!{nonenumerable}। অর্থাৎ, দেখাও যে $f_1$, $f_2$, \dots{} যদি অপেক্ষকগুলির একটি তালিকা হয় এবং প্রত্যেক $f_i\colon
\Nat \to \Nat$ হয়, তবে এমন কোনো $g \colon \Nat \to
\Nat$ আছে যা এই তালিকায় নেই।
\end{prob}

\end{document}
